\documentclass[12pt]{article}
\usepackage[utf8]{inputenc}
\usepackage[T1]{fontenc}
\usepackage{hyperref}
\title{Penan, Western Language Resources for K-12 Curriculum Development}
\author{Compiled by Claude Sonnet 5.5 (claude-sonnet-5-5) for Andr\'as Kornai}
\date{}
\begin{document}
\maketitle

\section{Introduction}
Western Penan (ISO 639-3 \texttt{pne}, Glottocode west2563) is an Austronesian language of Sarawak, Malaysia. It is distinct from Eastern Penan (\texttt{pez}). Speaker estimates vary: about 3,400 (ASJP), 3,800 (Global Recordings Network), 4,600 (Joshua Project). Scripture work exists, but the search found no Western Penan dictionary, grammar or school materials.

\section{Primary Educational Resources}
No primers or textbooks for Western Penan were found. A word list for Eastern Penan school children exists (see below) but concerns a different language.

\section{Linguistic and Reference Materials}
\begin{itemize}
\item Wikipedia, ``Penan language''. \url{https://en.wikipedia.org/wiki/Penan_language}
\item ASJP, Western Penan word list entry. \url{https://asjp.clld.org/languages/WESTERN_PENAN}
\item ELDP archive, ``A dictionary of Eastern Penan'' (P. Sercombe). This is Eastern Penan, not Western; it is listed only as a possible comparative resource. \url{https://eldp.access.preservica.com/dk0250/}
\end{itemize}

\section{Religious and Community Resources}
\begin{itemize}
\item Scripture Earth, Western Penan (\texttt{pne}) scripture resources. \url{https://www.scriptureearth.org/00eng.php?iso=pne}
\item Joshua Project, Penan-Nibong in Malaysia: lists a New Testament (1974) and portions (1958--1967); no complete Bible. \url{https://joshuaproject.net/people_groups/14482}
\item Digital Bible Society, New Testament in Eastern Penan (a different language, \texttt{pez}). \url{https://dbs.org/bibles/PEZBSM}
\end{itemize}

\section{Digital and Audio Resources}
\begin{itemize}
\item Global Recordings Network, Penan (\texttt{pne}) audio recordings. \url{https://globalrecordings.net/en/language/pne}
\end{itemize}

\section{Curriculum Development Context}
Western Penan is not known to be taught in schools. The New Testament and GRN audio are the only substantial texts found, and could seed a corpus and orthography study. Eastern Penan materials may be a partial model but are a separate language.

\end{document}